\documentclass[a4paper,11pt]{article}
\usepackage[utf8]{inputenc}
\usepackage[T1]{fontenc}
\usepackage[french]{babel}
\usepackage{graphicx}
\usepackage{amsmath}
\usepackage{geometry}
\usepackage{fancyhdr}
\usepackage{hyperref}
\usepackage[usenames,dvipsnames,svgnames,table]{xcolor}
\usepackage{enumerate}
\usepackage{booktabs}
\usepackage{tcolorbox}
\tcbuselibrary{listings,skins,breakable}
\usepackage{sourcecodepro}
\usepackage{FiraSans}
\usepackage{environ}
\usepackage{tikz}
\usepackage{pgfplots}
\pgfplotsset{compat=1.18}
\usetikzlibrary{arrows.meta,positioning,calc}

\geometry{a4paper, scale=0.8}
\urlstyle{same}

%% ─── Colors ──────────────────────────────────────────────────────
\definecolor{modernblue}{RGB}{66,133,244}
\definecolor{moderngray}{RGB}{95,99,104}
\definecolor{darkgray}{RGB}{60,64,67}
\definecolor{accentorange}{RGB}{255,111,0}
\definecolor{codegreen}{RGB}{15,157,88}
\definecolor{codeblue}{RGB}{25,118,210}
\definecolor{codebg}{RGB}{250,250,250}
\definecolor{codepurple}{RGB}{156,39,176}
\definecolor{niceblue}{RGB}{0,128,172}
 
%% ─── Code listing: PYTHON ───────────────────────────────────────
\newtcblisting{pythoncode}[1][]{
    listing engine=listings, colback=codebg, colframe=moderngray!50,
    listing only, left=5mm, enhanced, sharp corners=all, boxrule=0.5pt, top=0pt,
    listing options={
        language=Python, basicstyle=\ttfamily\scriptsize,
        keywordstyle=\color{codeblue}\bfseries,
        commentstyle=\color{codegreen}\itshape,
        stringstyle=\color{accentorange},
        numberstyle=\tiny\color{moderngray}, numbers=left, numbersep=8pt,
        breaklines=true, tabsize=4, showstringspaces=false,
        keepspaces=true, columns=flexible,
        literate=
            {é}{{\'{e}}}1 {è}{{\`{e}}}1 {ê}{{\^{e}}}1 {ë}{{\"{e}}}1
            {û}{{\^{u}}}1 {ù}{{\`{u}}}1 {î}{{\^{i}}}1 {ï}{{\"{i}}}1
            {ô}{{\^{o}}}1 {â}{{\^{a}}}1 {à}{{\`{a}}}1 {ä}{{\"{a}}}1
            {ç}{{\c{c}}}1 {Ç}{{\c{C}}}1
            {É}{{\'{E}}}1 {È}{{\`{E}}}1 {Ê}{{\^{E}}}1,
        #1
    }, breakable
}

%% ─── Lesson box ──────────────────────────────────────────────────
\newtcolorbox{lesson}[1][]{
    enhanced, breakable,
    colback=modernblue!5, colframe=modernblue!80,
    colbacktitle=modernblue, coltitle=white,
    fonttitle=\sffamily\bfseries, fontupper=\sffamily\small,
    boxrule=0.5pt, arc=3pt, left=8pt, right=8pt, top=8pt, bottom=8pt,
    attach boxed title to top left={yshift=-2mm, xshift=5mm},
    boxed title style={sharp corners, boxrule=0.5pt},
    title={#1}, before skip=1em, after skip=2em
}

%% ─── Solution toggle ─────────────────────────────────────────────
\newif\ifshowsolutions
\showsolutionsfalse

\ifshowsolutions
    \newtcolorbox{solution}{colback=white, colframe=white, breakable,
        fontupper=\color{red}}
\else
    \NewEnviron{solution}{}{}
\fi

%% ─── Question command ─────────────────────────────────────────────
\newcounter{question}[section]
\newcommand{\q}[1]{%
    \vspace{0.5em}\noindent
    {\color{black}\bfseries\stepcounter{question}\Alph{section}\arabic{question}.}
    \hspace{0.5em}#1\par\vspace{0.3em}
}
\renewcommand{\thesection}{{\Alph{section}}}

%% ─── Header / footer ─────────────────────────────────────────────
\pagestyle{fancy}
\fancyhf{}
\fancyhead[L]{2025-2026}
\fancyhead[C]{\textbf{\courseName}}
\fancyhead[R]{\large GEII -- BUT2}
\fancyfoot[C]{\thepage}
\fancyfoot[L]{\includegraphics[height=3em]{../../../common/logos_couleurs/IUTA.png}}
% \fancyfoot[R]{\href{mailto:prenom.nom@univ-smb.fr}{prenom.nom@univ-smb.fr}}

%% ─── Course metadata ─────────────────────────────────────────────
\newcommand{\courseName}{Outils Mathématiques et Logiciel --- OML3}
\newcommand{\chapterTitle}{TP 3 --- Transformée de Fourier}

\makeatletter
\renewcommand\maketitle{{\raggedright\begin{center}
    {\Huge\bfseries\@title}
\end{center}}}
\makeatother

\usepackage{amsmath,amsfonts,bm}

%% ─── Operators ───────────────────────────────────────────────────
\DeclareMathOperator*{\argmax}{arg\,max}
\DeclareMathOperator*{\argmin}{arg\,min}
\DeclareMathOperator{\MSE}{MSE}
\DeclareMathOperator{\sign}{sign}
\DeclareMathOperator{\sinc}{sinc}
\DeclareMathOperator{\tr}{tr}

%% ─── Number sets ─────────────────────────────────────────────────
\def\R{{\mathbb{R}}}
\def\C{{\mathbb{C}}}
\def\N{{\mathbb{N}}}
\def\Z{{\mathbb{Z}}}

%% ─── Infinitesimal and units ────────────────────────────────────
\def\d{{\rm d}}
\def\Hz{{\mbox{Hz}}}
\def\m{{\mbox{m}}}

%% ─── Complex / signal ───────────────────────────────────────────
\DeclareMathOperator{\complexj}{\mathsf{j}}   % imaginary unit

%% ─── Probability ─────────────────────────────────────────────────
\newcommand{\Esp}[2]{\mathbb{E}_{#1}\!\left[#2\right]}
\renewcommand{\Pr}{\mathbb{P}}

%% ─── Linear algebra ──────────────────────────────────────────────
\newcommand{\abs}[1]{\left|{#1}\right|}
\newcommand{\norm}[1]{\left\|#1\right\|}
\newcommand{\inner}[2]{\left\langle#1,\,#2\right\rangle}

%% ─── Logic ───────────────────────────────────────────────────────
\newcommand{\then}{\Rightarrow}
\newcommand{\iif}{\Leftrightarrow}

%% ─── Draft annotation ────────────────────────────────────────────
\newcommand{\red}[1]{{\color{red}#1}}

%% ─── Shared figure color (used by common/figures/*/*.tex bodies) ──
\definecolor{niceblue}{RGB}{0,128,172}

%% ─── Rappel header (styled subtitle for reference-card asides) ───
%% Requires \definecolor{niceblue}{...} to be defined by the including document.
\newcommand{\rappelhead}[1]{%
    \bigskip
    {\color{niceblue}\bfseries #1}\\[-0.4em]
    {\color{niceblue!35}\rule{\linewidth}{0.6pt}}
    \smallskip

}


\title{\chapterTitle}
\author{}
\date{}

%% ─────────────────────────────────────────────────────────────────
\begin{document}
\fontfamily{bch}\selectfont
\fontsize{12pt}{14pt}\selectfont

\maketitle
\thispagestyle{fancy}


%% ─── AVANT DE COMMENCER ──────────────────────────────────────────
\section*{Avant de commencer}

Ce TP est écrit en Python. Les fonctions dont vous aurez besoin sont déjà
écrites. Votre travail consiste à exécuter les cellules en modifiant ou
non les paramètres, et à lire les tracés produits.

\textbf{Ce document est le sujet de référence.} Le notebook
\texttt{TP3.ipynb} est un compagnon d'exécution : il ne reprend que
l'essentiel de chaque question, repérée par son numéro (A1, A2, B1,
\dots). Gardez ce PDF ouvert pendant tout le TP.

\textbf{Vos réponses vont sur feuille.}

\subsection*{Configuration de l'environnement}
\begin{enumerate}
    \item Ouvrir \texttt{TP3.ipynb} dans VsCode.
    \item Exécuter la première cellule (import des fonctions fournies,
          dans \texttt{tp3\_helpers.py}) : rien ne devrait s'afficher,
          c'est normal.
    \item Exécuter ensuite les cellules dans l'ordre, en vous référant aux
          questions ci-dessous.
\end{enumerate}

\vspace{2em}\hrule\vspace{2em}
\newpage


%% ─── SECTION A ──────────────────────────────────────────────────
\section{Valeur moyenne et exponentielle complexe}

Soit $\Pi(t)$ la fonction porte de largeur $1$, hauteur $1$, centré en $0$.

\q{Rappeler la définition de la transformée de Fourier appliqué à la fonction porte.}
\begin{solution}
    $$X(f) = \int_{-\infty}^{+\infty} \Pi(t) e^{-\complexj 2\pi f t}$$
\end{solution}


\q{Calculer $X(f)$, la transformée de Fourier de la fonction porte.}
\begin{solution}
    $\displaystyle\int_{-1/2}^{1/2} e^{-2j\pi f_0 t}\,\mathrm dt
    = \frac{\sin(\pi f_0)}{\pi f_0} = \sinc(\pi f_0)$.
\end{solution}

\begin{lesson}[Fonctions \texttt{fourier\_transform}, \texttt{explorer}]
    \texttt{fourier\_transform(signal, f0)} calcule cette même intégrale
    par intégration numérique, pour un \texttt{Signal} et une fréquence
    \texttt{f0} quelconques. \texttt{explorer(f0)} affiche, pour la porte,
    le produit avec $\cos(2\pi f_0t)$ (et avec $\sin(2\pi f_0t)$ si
    \texttt{montrer\_sin=True}), l'aire signée sous chaque produit, sa
    moyenne, et compare celle-ci à la valeur théorique de la question A1.
\end{lesson}

\q{Exécuter la cellule \texttt{explorer(f0=1)}. Elle affiche, pour la
porte, le produit avec $\cos(2\pi f_0t)$ et avec $\sin(2\pi f_0t)$,
$$
    A\Pi_\theta(t) \times \cos(2\pi f_0 t) \quad \text{et} \quad A\Pi_\theta(t) \times \sin(2\pi f_0 t).
$$
L'aire est coloriée (positive en bleu et négative en rouge), la valeur moyenne (qui correspond à l'aire sous la courbe) apparait en pointillés.
\textbf{Identifier ces deux moyennes comme la partie réelle et (l'opposé de) la partie imaginaire de} $X(f_0)$. \textbf{Pourquoi la moyenne du produit en sinus est-elle nulle ici ?}
}
\begin{solution}
    Par définition,
    $X(f_0)=\int f(t)e^{-2j\pi f_0t}\,\mathrm dt = \int f(t)\cos(2\pi f_0t)\,\mathrm dt - j\int f(t)\sin(2\pi f_0t)\,\mathrm dt$ :
    les deux moyennes affichées sont exactement $\mathrm{Re}(X(f_0))$ et
    $-\mathrm{Im}(X(f_0))$. La moyenne en sinus est nulle ici car la
    porte est paire et $\sin$ est impair : leur produit est impair,
    d'intégrale nulle sur un intervalle symétrique.
\end{solution}

\q{Reprendre \texttt{explorer} (partie en cosinus seule) pour
$$f_0=0.1, f_0=1, f_0=1.43 ,f_0=2.$$ Comparer, à chaque fois,
la valeur théorique (formule de A2) et la valeur numérique. Que se
passe-t-il en $f_0=1$ et $f_0=2$ ? Que remarquez-vous sur le signe de la
moyenne en $f_0=1.43$, comparé à $f_0=0.1$ ?}
\begin{solution}
    Les valeurs théorique et numérique coïncident à chaque fois (aux
    erreurs d'arrondi près), confirmant le calcul de A2. En $f_0=1$ et
    $f_0=2$, la moyenne est nulle : la porte contient alors un nombre
    entier de périodes du cosinus, et les lobes positifs et négatifs du
    produit se compensent exactement. En $f_0=1.43$, la moyenne est non
    nulle mais de signe opposé à celle de $f_0=0.1$ : on a franchi un
    zéro et changé de lobe. Ces quatre valeurs sont respectivement dans
    le lobe principal, sur le premier zéro, sur le maximum du premier
    lobe secondaire, et sur le second zéro de $X(f)$ -- leur observation
    trace déjà la forme du spectre étudié formellement en B.
\end{solution}

%% ─── SECTION B ──────────────────────────────────────────────────
\section{La transformée de Fourier}

\q{Que devient la transformée de Fourier de la porte pour une largeur $\theta$ quelconque, une hauteur de 1 et centrée en 0 ? 
On notera la fonction $\Pi_\theta(t)$.}
\begin{solution}
    $X(f)=\displaystyle\int_{-\infty}^{+\infty}x(t)\,e^{-2j\pi ft}\,\mathrm dt$.
    Pour $\Pi_\theta$ :
    $X(f) = \theta\,\sinc(\pi f\theta)$, réelle (et paire), car $\Pi_\theta$
    est réelle et paire.
\end{solution}

\begin{lesson}[Fonctions \texttt{spectrum}, \texttt{plot\_spectrum}]
    \texttt{spectrum(signal, freqs)} calcule $X(f)$ (tableau complexe)
    pour chaque fréquence de \texttt{freqs}. \texttt{plot\_spectrum(freqs,
    X)} trace $|X(f)|$ (ou sa partie réelle/imaginaire avec
    \texttt{kind="re"}/\texttt{"im"}).
\end{lesson}

\q{Tracer, superposées, les transformées de portes de plusieurs largeurs
$\theta$. Comment la largeur du lobe principal évolue-t-elle avec
$\theta$ ? Où le spectre s'annule-t-il ?}
\begin{solution}
    Le lobe principal a une largeur $\propto1/\theta$ : plus la porte est
    étroite (temps), plus son spectre est étalé (fréquence), et
    inversement. $X(f)=0 \iff \sin(\pi f\theta)=0$ et $f\neq0$, soit
    $f=k/\theta$, $k\in\Z^*$.
\end{solution}

\q{Rétrécir $\theta$ tout en gardant l'aire de la porte égale à $1$
en fixant la hauteur à $1/\theta$. Vers quel objet le signal temporel converge-t-il ?
Et son spectre ?}
\begin{solution}
    Le signal temporel converge vers l'impulsion de Dirac $\delta(t)$
    (aire $1$, largeur nulle). Son spectre,
    $\frac1\theta\cdot\theta\sinc(\pi f\theta)=\sinc(\pi f\theta)$,
    converge vers la constante $1$ pour tout $f$.
\end{solution}

\q{Quel est le support (l'ensemble des $t$ où le signal est non nul) de
de la fonction porte ? Quel est le support de son spectre ? Comment évolue-t-il quand la largeur de la porte rétrécie et inversement ?
Qu'en déduisez-vous, plus généralement, sur le lien entre support temporel fini et support
fréquentiel ?}
\begin{solution}
    Support de $\delta$ : $\{0\}$, réduit à un point. Support de son
    spectre (constant, non nul) : $\R$ tout entier. Un signal
    concentré dans le temps a nécessairement un spectre étalé en
    fréquence, et réciproquement : un support temporel fini interdit un
    support fréquentiel fini (sauf le signal nul).
\end{solution}

\q{On tronque le spectre de $\Pi_\theta$ en le multipliant par une porte
fréquentielle de largeur $B$.
Exprimer $X_{\text{tronqué}}(f)$ à partir de $X(f)$.

\emph{Indice : $X_{\text{tronqué}}$ ne
peut être non nul que là où la porte fréquentielle vaut $1$.}}
\begin{solution}
    $X_{\text{tronqué}}(f) = X(f)$ pour $|f|<B/2$, $0$ sinon, avec
    $X(f)=\theta\sinc(\pi f\theta)$ (B1).
\end{solution}
\q{Exprimer $x_{\text{tronqué}}(t)$ le signal reconstruit comme étant le produit de convolution de deux signaux.}
\begin{solution}
    Le produit en fréquentiel devient convolution en temporel, d'où 
    $$
        x_{\textbf{tronqué}}(t) = x(t) * k(t) 
    $$
    où $k(t)$ est la transformée de Fourier inverse de la porte fréquentielle qui par dualité vaut 
    $$
        k(t) = B\sinc(\pi t B)
    $$
\end{solution}

\begin{lesson}[Fonction \texttt{inverse\_fourier\_transform}]
    \texttt{inverse\_fourier\_transform(X, t, fmin, fmax)} calcule
    $\displaystyle\int_{f_{\min}}^{f_{\max}} X(f)\,e^{2j\pi ft}\,\mathrm df$
    par intégration numérique, pour \texttt{X} une fonction $f\mapsto X(f)$
    quelconque.
\end{lesson}


\q{Executer la cellule de reconstruction et observé $x_{\text{tronqué}}(t)$ pour plusieurs largeurs de bande $B$. Comment s'appelle le phénomène que l'on observe aux bords ? }
\begin{solution}
    Un dépassement apparaît de part et d'autre de chaque bord et ne
    disparaît pas quand $W$ augmente : il se resserre autour du bord sans
    que son amplitude relative ne diminue -- le phénomène de Gibbs, ici en
    temps, causé par une troncature en fréquence (au lieu d'une
    troncature en temps). Ailleurs, la reconstruction se rapproche de la
    porte d'origine.
\end{solution}


%% ─── SECTION C ──────────────────────────────────────────────────
\section{Résolution fréquentielle}

On observe, sur une durée finie $T$, la somme de deux signaux
sinusoïdaux de fréquences proches $f_1$ et $f_2$ :
$$x(t) = \Pi_{T}(t)\left[\cos(2\pi f_1 t) + \cos(2\pi f_2 t)\right]$$
(C'est un battement, comme lorsqu'on accorde deux instruments.)

\q{Afficher $x(t)$ pour $T=3$. Les fréquences $f_1$ et $f_2$
sont-elles reconnaissables sur ce tracé ?}
\begin{solution}
    Le tracé ressemble à une simple sinusoïde unique, de fréquence
    intermédiaire : sur une durée aussi courte, rien ne distingue à
    l'œil la somme de deux fréquences proches d'une fréquence unique.
\end{solution}

\q{Rappeler l'expression de la transformée de Fourier de $\cos(2\pi f_0 t)$ et en déduire la transformée de Fourier de $x(t)$, $X(f)$. }
\begin{solution}
    $\mathcal F[\cos(2\pi f_0t)](f) = \frac12\delta(f-f_0)+\frac12\delta(f+f_0)$
    (rappel du cours). Par linéarité,
    $$
        \mathcal F[\cos(2\pi f_1t)+\cos(2\pi f_2t)](f) =
        \tfrac12\big[\delta(f-f_1)+\delta(f+f_1)+\delta(f-f_2)+\delta(f+f_2)\big].
    $$
    $x(t)$ est ce signal multiplié par la porte temporelle $\Pi_T(t)$ :
    comme en B6 mais dans l'autre sens (c'est ici le temps qui est
    tronqué), le produit en temps devient une convolution en fréquence
    par le spectre de $\Pi_T$, soit $T\sinc(\pi fT)$ :
    $$
        X(f) = \tfrac12\big[\delta(f-f_1)+\delta(f+f_1)+\delta(f-f_2)+\delta(f+f_2)\big] * T\sinc(\pi fT).
    $$
    Convoluer avec un Dirac translate : $X(f)$ est la somme de quatre
    lobes en $\sinc$, centrés en $\pm f_1,\pm f_2$, chacun de largeur
    $\propto1/T$ -- exactement ce qui est observé numériquement en
    C1--C4.
\end{solution}

\q{Calculer et tracer le spectre pour ce même $T$. Où se situe le pic ?
Y a-t-il un seul pic ?}
\begin{solution}
    Un seul pic apparaît, situé entre $f_1$ et $f_2$ sans coïncider avec
    aucun des deux : comme en B2/B3, chaque fréquence contribue un lobe
    spectral de largeur $\propto 1/T$ ; ici $T$ est trop court, les deux
    lobes se recouvrent et fusionnent en un seul pic.
\end{solution}

\q{Reprendre C1 et C3 avec $T=8$. Que devient le tracé temporel ? Et le
spectre ?}
\begin{solution}
    Dans le temps, une enveloppe apparaît maintenant clairement (le
    battement, de période $1/|f_1-f_2|$) -- mais $f_1$ et $f_2$ restent
    impossibles à lire précisément sur ce tracé. Dans le spectre en
    revanche, les deux lobes se sont assez rétrécis pour se séparer :
    deux pics distincts apparaissent, chacun proche de $f_1$ et $f_2$.
\end{solution}

\q{Comparer le spectre pour plusieurs durées $T$ (cellule suivante). À
partir de quelle durée les deux pics se séparent-ils clairement ?}
\begin{solution}
    Les deux pics restent fusionnés pour $T=3$ et $T=5$, et sont
    nettement séparés dès $T=8$ : la séparation a lieu quelque part
    entre $T=5$ et $T=8$.
\end{solution}

\q{Sur papier, à partir de la largeur des lobes observée en B2, à
quelle condition sur $T$ peut-on espérer séparer deux pics distants de
$|f_1-f_2|$ ?}
\begin{solution}
    D'après B2, chaque lobe a une largeur $\propto 1/T$. Séparer deux
    pics distants de $|f_1-f_2|$ demande que cette largeur devienne
    plus petite que $|f_1-f_2|$, soit $T\gtrsim 1/|f_1-f_2|$ -- ici
    $|f_1-f_2|=0.2$, donc $T\gtrsim5$, cohérent avec l'observation de
    C5.
\end{solution}

\end{document}
